\chapter{Finite-Depth Adaptive Formation at One Base Point}
\label{ch:finite-depth-formation}

The normalized parent construction and the memory results of
\Cref{ch:agent-network-rg,ch:inference-memory} distinguish an exact
probability reduction from an autonomous reduced dynamics. This chapter
constructs a finite family in which successive reductions can actually be
used during inference. The family has independent moving interaction
links, retained network loops, a nonlinear observation of source-dependent
belief mismatch, and positive-rate changes of the actual source
assignments. Every new parent begins with the preceding actual retained
law and subsequently learns from its own law. The quantitative result is
\Cref{thm:finite-depth-adaptive-formation}; its estimates are proved in
\Cref{app:finite-depth-formation-proofs}. \status{ESTABLISHED}

The construction adopts one fixed base point \(c_*\), a finite
triangle-chain graph, supplied interaction-energy contrasts, a proper
correlated reference, a compatible observation, and calibration-assisted
learning of one scalar model parameter. It is a realization inside the
principal/associated-bundle framework, with no base transport. The
continuous mobility, categorical proposal rule, and model metric are
declared inference kinetics. A scalar VFE alone does not choose them.
\status{HYPOTHESIS}

\section{The pointwise carrier and its restoring flag}
\label{sec:rg-formation-carrier}

Fix an integer \(L\ge2\), set \(b=2^{L-1}\), and arrange \(b\) disjoint
triangles in a chain. A bridge joins each consecutive pair of triangles.
Each triangle is cyclically oriented; every other edge has one fixed
orientation. The graph has \(n=3b\) vertices and \(4b-1\) edges. Bridge
number \(j\), between triangles \(j\) and \(j+1\), has level
\(\ell(j)=2+v_2(j)\), where \(v_2(j)\) is the exponent of two dividing
\(j\). Triangle edges have level one. The numbers \(J_1\ge\cdots\ge
J_L>0\) introduced below are energy strengths, not categorical source
probabilities. \status{DEFINITION}

Use the translation action of \((\R,+)\) on two scalar law fibers at
\(c_*\), and fix \(0<c<1\). A configuration and its law presentations are
\begin{equation}
\begin{aligned}
x&=(u,v,\zeta)\in\R^{2n+4b-1}=\R^{10b-1},\\
q_i^b(\cdot\mid x)&=\mathcal N(u_i,1),&
q_i^m(\cdot\mid x)&=\mathcal N(v_i,c),\\
\Theta_e(k)&=k+\zeta_e,&
d_i(x)&=u_i-v_i,\\
r_e(x)&=u_i-u_j-\zeta_e\quad(e:j\to i).
\end{aligned}
\label{eq:rg-formation-configuration}
\end{equation}
Here \(x\) is a random configuration of law presentations; it is not the
fixed structural datum \(X\) of the general construction. The same local
translation acts on both fibers in this specialization. The independent
coordinates \(\zeta_e\) are graph links, with
\(\zeta_{\bar e}=-\zeta_e\). No inter-agent model-alignment record is
included. This choice leaves the local model evaluator and model learning
present. \status{DEFINITION}

Under a passive vertex-frame translation \(a_i\), the means become
\(u_i+a_i,v_i+a_i\) and an edge coordinate becomes
\(\zeta_e+a_i-a_j\). Thus \(d_i\) and \(r_e\) are invariant. Reference
centers and observation records transform with the same passive change;
holding those centers fixed while moving the configuration would instead
change the model. The proper reference used below is not an invariant
Haar probability on a noncompact group. These are the coordinate and
reference conventions of this specialization. \status{DEFINITION}

For triangle \(q\), let \(\mathfrak h_q=\sum_{e\in q}\zeta_e\) be its
oriented translation holonomy and put \(h_q=\mathfrak h_q/\sqrt3\).
The normalized loop vector is \(h=(h_q)_{q=1}^b\). With
\begin{equation}
Q_\triangle=
\begin{pmatrix}
1/\sqrt2&-1/\sqrt2&0\\
1/\sqrt6&1/\sqrt6&-2/\sqrt6
\end{pmatrix},
\qquad
|r_q|^2=|Q_\triangle r_q|^2+h_q^2,
\label{eq:rg-formation-loop-decomposition}
\end{equation}
the second identity follows from
\(\sum_{e\in q}r_e=-\mathfrak h_q\). In particular the loop energy
in this decomposition belongs to the original sum of edge mismatches.
It is not an additional flattening term. \status{ESTABLISHED}

Let \(B_1\) stack the \(n\) rows \(d_i\) and the \(2b\) internal
contrast rows \(Q_\triangle r_q\). For \(s\ge2\), let \(B_s\) stack
the bridge residuals of level \(s\), and write \(B\) for their stack.
The link-coordinate columns in a linear dependence among these rows force
all contrast and bridge coefficients to vanish. The model-mean columns
then force the remaining coefficients to vanish. Hence
\(\operatorname{rank}B=6b-1\). Every normalized loop covector is
orthogonal to every row of \(B\). If \(D\) is the oriented graph
incidence matrix, define
\begin{equation}
\begin{aligned}
\mathcal K&=\{(t,t,Dt):t\in\R^n\},&
\mathcal H&=\operatorname{span}\{\text{normalized loop covectors}\},\\
\mathcal N&=\operatorname{range}B^\top,&
\R^{10b-1}&=\mathcal N\oplus\mathcal H\oplus\mathcal K,\\
H&=B^\top\operatorname{diag}(J_sI)B+J_1P_{\mathcal H}.
\end{aligned}
\label{eq:rg-formation-orthogonal-split}
\end{equation}
The direct sum is orthogonal. The matrix \(H\) has kernel
\(\mathcal K\), and
\(x^\top Hx=J_1\sum_i d_i(x)^2+\sum_eJ_{\ell(e)}r_e(x)^2\).
These are established finite-dimensional identities.
\status{ESTABLISHED}

Define the nested retained spaces and projections by
\begin{equation}
\begin{aligned}
W_s&=\ker(B_1,\ldots,B_s),&
\mathcal N_s&=W_s\cap\mathcal N,&
M_s&=\mathcal N_s\oplus\mathcal H,\\
P_s&:\R^{10b-1}\longrightarrow W_s,&
P_0&=I,&
P_sP_p&=P_s\quad(p<s).
\end{aligned}
\label{eq:rg-formation-projections}
\end{equation}
Each \(P_s\) is the Euclidean orthogonal projection. All loop coordinates
and all of \(\mathcal K\) survive. The final space is
\(W_L=\mathcal H\oplus\mathcal K\). Write \(k_p=\dim M_p\) and
\(d_p=\dim W_p^\perp\); the latter is the cumulative eliminated
dimension. The elementary row count gives
\begin{equation}
\dim W_s=4b-1+\frac{b}{2^{s-1}}\quad(1\le s\le L),
\qquad
\dim W_L=4b.
\label{eq:rg-formation-dimensions}
\end{equation}
These dimensions count retained continuous coordinates rather than
agents. \status{ESTABLISHED}

Choose orthonormal columns \(\mathsf Q_s\) spanning
\(\mathcal N_{s-1}\cap\mathcal N_s^\perp\), and write
\(x_{\mathcal N}=\sum_s\mathsf Q_sw_s\). The matrix
\(\mathsf B=B[\mathsf Q_1\cdots\mathsf Q_L]\), including its inverse, is block lower
triangular. Conjugation by
\(\operatorname{diag}(\sqrt{J_s}I)\) cannot enlarge a permitted lower
block by a factor greater than one. Therefore
\begin{equation}
\|\mathsf B^{-1}\|_F^{-2}\sum_sJ_s|w_s|^2
\le x_{\mathcal N}^{\top}Hx_{\mathcal N}
\le\|\mathsf B\|_F^2\sum_sJ_s|w_s|^2.
\label{eq:rg-formation-weighted-flag}
\end{equation}
Minimizing this quadratic form over an initial group of normal layers
gives the same bounds for its Schur complement. These explicit finite
matrix constants underlie the restoring estimates in the proof appendix.
\status{ESTABLISHED}

\section{One normalized experiment with a nonlinear source record}
\label{sec:rg-formation-normalized-model}

At each vertex the model sample is \(m_i\sim\mathcal N(v_i,c)\), and
the evaluator sends \(m_i\) to \(\mathcal N(m_i,1-c)\) on the belief
sample \(k_i\). Among couplings with the prescribed marginals
\(\mathcal N(v_i,c)\) and \(\mathcal N(u_i,1)\), define
\begin{equation}
m_i=v_i+\xi_i,\qquad
k_i=u_i+\xi_i+\varepsilon_i,\qquad
\xi_i\sim\mathcal N(0,c),\quad
\varepsilon_i\sim\mathcal N(0,1-c),
\label{eq:rg-formation-evaluator-lift}
\end{equation}
where the two Gaussian noises are independent. This is the chosen
conditional recognition coupling at a fixed configuration.
\status{DEFINITION}

Its log-density ratio to the physical model/evaluator coupling is
\(\{d_i(k_i-m_i)-d_i^2/2\}/(1-c)\). Its expectation is fixed by
the two prescribed marginals. Thus, for any other admissible coupling
\(\nu_i\), subtracting this ratio gives
\begin{equation}
\KL(\nu_i\Vert
  \mathcal N(v_i,c)\operatorname{ev}_i)
=\KL(\nu_i\Vert\nu_i^*)+\frac{d_i^2}{2(1-c)}.
\label{eq:rg-formation-evaluator-optimum}
\end{equation}
Here \(\nu_i^*\) is the law in
\cref{eq:rg-formation-evaluator-lift}. This proves its optimality and
the total restriction cost
\(d^{\mathrm{ev}}(x)=x^\top Rx/2\), where
\(R=b_d^\top b_d/(1-c)\) and \(b_dx=u-v\).
\status{ESTABLISHED}

Let \(\vartheta\) be the learned scalar parameter. The physical
configuration prior and its Lebesgue normalizer are
\begin{equation}
\begin{aligned}
A_0(\vartheta)&=\vartheta H+P_{\mathcal K}-TR,\\
\mathbb P^0_\vartheta(dx)
&=Z_0(\vartheta)^{-1}
  \exp\left[-\frac{x^\top A_0(\vartheta)x}{2T}\right]dx,\\
Z_0(\vartheta)&=(2\pi T)^{(10b-1)/2}
  \det A_0(\vartheta)^{-1/2}.
\end{aligned}
\label{eq:rg-formation-physical-prior}
\end{equation}
Assume \(T>0\), \(1\le\vartheta\le21/10\), and
\(T\le J_1(1-c)/2\). Then \(TR\preceq H/2\), so
\(A_0(\vartheta)\succeq(\vartheta/2)H+P_{\mathcal K}\succ0\).
The prior is a normalized correlated Gaussian. Its projector
\(P_{\mathcal K}\) is generally dense; no incident-edge-only force
evaluation is claimed. \status{ESTABLISHED}

Fix \(0<\kappa\le1\) and a nonzero \(z_\star\in\mathcal K\) whose
node entries are not all equal. Declare the observation kernels
\begin{equation}
O\mid x\sim
  \mathcal N_{\mathcal K}(P_{\mathcal K}x,(T/\kappa)I_{\mathcal K}),
\qquad
Y_{\mathrm{cal}}\mid\vartheta\sim\mathcal N(\vartheta,T),
\label{eq:rg-formation-observation-kernels}
\end{equation}
and retain the records \(O=z_\star\), \(Y_{\mathrm{cal}}=2\).
The first Gaussian is a density on the \(n\)-dimensional carrier
\(\mathcal K\), not a singular full-space density. Its compatibility
with the coherent root and the separate calibration record are modeling
assumptions. \status{HYPOTHESIS}

Each receiver \(i\) has the two candidate sources
\(\mathcal N_i^\triangle\) inside its own triangle. The complete
assignment is \(S=(S_i)_i\in\mathcal J=\prod_i\mathcal N_i^\triangle\),
with \(|\mathcal J|=2^n\). Every bridge retains its original interaction
energy but is excluded from these source palettes. Choose one
distinguished receiver \(i_0\) and label its two sources \(+\) and
\(-\). The independent source priors are
\begin{equation}
\begin{aligned}
\pi_{i_0,+}(\vartheta)
  &=\operatorname{sigmoid}(4(\vartheta-3/2)),&
\pi_{i_0,-}(\vartheta)&=1-\pi_{i_0,+}(\vartheta),\\
\pi_{i,j}(\vartheta)&=\tfrac12\quad(i\ne i_0),&
\pi_S(\vartheta)&=\prod_i\pi_{i,S_i}(\vartheta).
\end{aligned}
\label{eq:rg-formation-source-priors}
\end{equation}
This internal-neighbor support and the supplied prior parameterization
are adopted restrictions. \status{HYPOTHESIS}

Set \(B_0=10^{-4}\), \(\delta=B_0/n\), and
\(f(r)=1-e^{-r^2/2}\). At each receiver add a binary observation
\(E_i\), conditionally independent given \((x,S)\), with
\begin{equation}
\begin{aligned}
\Pr(E_i=1\mid x,S_i=j)
  &=q_0\exp\left[-\delta
       f\left(\sqrt{J_1/T}r_{ij}(x)\right)\right],
       \qquad q_0=\tfrac12,\\
\phi_S(x)&=\delta\sum_i
       f\left(\sqrt{J_1/T}r_{i,S_i}(x)\right),
       \qquad 0\le\phi_S\le B_0.
\end{aligned}
\label{eq:rg-formation-source-record}
\end{equation}
The complementary probability defines \(E_i=0\), and the observed
record is \(E_i=1\) for every \(i\). The standardized square satisfies
\(J_1r_{ij}^2/(2T)=(J_1/T)
\KL(\mathcal N(u_i,1)\Vert\mathcal N(u_j+\zeta_{ij},1))\).
The bounded function \(f\) specifies an additional normalized
observation of that mismatch; it is not identified with the mismatch
itself. The coefficient \(\delta\) is fixed at every prescribed depth
and is not taken to zero in the theorem. \status{DEFINITION}

Multiplying the normalized prior, the two observation kernels, the
source priors, the Bernoulli kernels, and the model/evaluator kernels
defines one complete generative joint:
\begin{equation}
\begin{aligned}
\mathbb P_\vartheta
={}&\mathbb P^0_\vartheta(dx)
 L_O(dO\mid x)L_{\mathrm{cal}}(dY_{\mathrm{cal}}\mid\vartheta)\\
&{}\times\prod_i\left[
 \pi_i(dS_i\mid\vartheta)L_i(dE_i\mid x,S_i)
 \mathcal N(dm_i;v_i,c)
 \mathcal N(dk_i;m_i,1-c)\right].
\end{aligned}
\label{eq:rg-formation-full-joint}
\end{equation}
Iterated integration of its conditional factors gives total mass one.
None of its kernels takes a recognition law as an input. The selected
observation has a strictly positive density. This establishes the
normalized experiment used throughout the chapter.
\status{ESTABLISHED}

\section{Physical posterior, restricted optimizer, and actual recognition}
\label{sec:rg-formation-laws}

Write \(\mathsf o=(z_\star,2,\mathbf1)\). The chapter-local actual
outer law has a smooth positive density \(q(x,S)\) relative to
Lebesgue measure times counting measure. Its conditional evaluator
coupling is \(\prod_i\nu_i^*\) from
\cref{eq:rg-formation-evaluator-lift}. The outer density is otherwise
unrestricted: assignments and configurations may be correlated. It is
not required to use the source posterior as its conditional assignment
law. This is the recognition family on which the restricted VFE is
defined. \status{DEFINITION}

Set
\begin{equation}
\begin{aligned}
\widehat K_\vartheta&=\vartheta H+(1+\kappa)P_{\mathcal K},&
G_\vartheta(x)&=\tfrac12x^\top\widehat K_\vartheta x
                  -\kappa z_\star^\top x,\\
G_\vartheta^{\mathrm{phys}}(x)
   &=G_\vartheta(x)-\tfrac T2x^\top Rx,&
U_{\vartheta,S}(x)&=G_\vartheta(x)+T\phi_S(x),\\
Z_{\mathrm{opt}}(\vartheta)
   &=\sum_S\pi_S(\vartheta)
      \int e^{-G_\vartheta(x)/T-\phi_S(x)}dx,&
Z_{\mathrm{phys}}(\vartheta)
   &=\sum_S\pi_S(\vartheta)
      \int e^{-G_\vartheta^{\mathrm{phys}}(x)/T-\phi_S(x)}dx.
\end{aligned}
\label{eq:rg-formation-partition-integrals}
\end{equation}
Gaussian domination and \(0\le\phi_S\le B_0\) make these integrals
finite and positive. The local symbol \(\pi_\vartheta\) denotes the
normalized restricted optimizer
\(\pi_S(\vartheta)e^{-U_{\vartheta,S}(x)/T}/Z_{\mathrm{opt}}(\vartheta)\).
The physical posterior \(\boldsymbol\Pi_{\vartheta,\mathsf o}\) on
\((x,S)\) instead uses \(G_\vartheta^{\mathrm{phys}}\) and
\(Z_{\mathrm{phys}}\). Both have full support, and their conditional
assignment row at fixed \(x\) is
\begin{equation}
b_S(x,\vartheta)=
\frac{\pi_S(\vartheta)e^{-\phi_S(x)}}
     {\sum_{S'}\pi_{S'}(\vartheta)e^{-\phi_{S'}(x)}}.
\label{eq:rg-formation-full-source-row}
\end{equation}
It is the product of the receiver posterior rows. Equality of these
conditional source rows does not identify the two full laws or their
hidden-coordinate conditionals. \status{ESTABLISHED}

For explicit scalar accounting, define
\begin{equation}
\begin{aligned}
c_O&=(\kappa/(2\pi T))^{n/2},&
c_{\mathrm{cal}}&=(2\pi T)^{-1/2},\\
B_T(\vartheta)
 &=T\log Z_0(\vartheta)-T\log(c_Oc_{\mathrm{cal}}q_0^n)
   +\tfrac12\kappa|z_\star|^2+\tfrac12(\vartheta-2)^2,\\
C_T(\vartheta)&=B_T(\vartheta)-T\log Z_{\mathrm{opt}}(\vartheta),&
p_\vartheta(\mathsf o)&=e^{-B_T(\vartheta)/T}
                         Z_{\mathrm{phys}}(\vartheta).
\end{aligned}
\label{eq:rg-formation-evidence-scalar}
\end{equation}
Expansion of the complete joint and
\cref{eq:rg-formation-evaluator-optimum} gives
\begin{equation}
\begin{aligned}
F_T(q,\vartheta)
 &=T\sum_S\int q(x,S)\log q(x,S)dx\\
 &\quad+\E_q[
       G_\vartheta(x)+T\phi_S(x)-T\log\pi_S(\vartheta)]
       +B_T(\vartheta)\\
 &=T\KL(q\Vert\pi_\vartheta)+C_T(\vartheta).
\end{aligned}
\label{eq:rg-formation-restricted-vfe}
\end{equation}
The scalar is the scaled negative ELBO of the declared evaluator
restriction. Its difference from the unconstrained physical optimum is
\(C_T+T\log p_\vartheta(\mathsf o)
=T\log(Z_{\mathrm{phys}}/Z_{\mathrm{opt}})\ge0\).
The ratio, including its parameter derivative, is retained. These are
exact identities, not Gaussian approximations to the nonlinear law.
\status{ESTABLISHED}

For clarity, the actual conditional assignment density is
\(\rho_q(S\mid x)=q(x,S)/\sum_{S'}q(x,S')\). Its receiver row is
\(\beta^Q_{ij}(x)=\sum_{S:S_i=j}\rho_q(S\mid x)\); the model posterior
row \(\beta^P_{ij}(x,\vartheta)\) is obtained from
\cref{eq:rg-formation-full-source-row}. Receiver occupancy may be fixed
at \(1/n\) by a readout that selects a receiver and reports its current
source. The resulting event law is the pushforward of the complete
marked input law. It is not a product of independently averaged
occupancies, source rows, and memberships. These are the event and
recognition conventions of this family. \status{DEFINITION}

\section{Exact marked parents and their inherited inference}
\label{sec:rg-formation-parents}

The parent carrier at level \(p\) is \(W_p\times\mathcal J\).
For \(p<s\), define the normalized deterministic Markov kernel
\begin{equation}
C_{s\leftarrow p}((y,S),\cdot)=\delta_{(P_sy,S)},
\qquad C_s=C_{s\leftarrow0}.
\label{eq:rg-formation-marked-channel}
\end{equation}
The matrices \(P_s\) and the probability kernels \(C_{s\leftarrow p}\)
are different typed objects. Their composition follows from
\(P_sP_p=P_s\); no assignment is erased. The same channel, extended
by the identity on observations, pushes the generative joint, physical
posterior, restricted optimizer, and actual recognition. This is a
declared pointwise coarse construction. \status{DEFINITION}

Let \(\pi_\vartheta^p=(C_p)_\#\pi_\vartheta\), and similarly denote
the physical and actual pushes by
\(\boldsymbol\Pi_{\vartheta,\mathsf o}^p\) and \((C_p)_\#q\).
Full evaluator and boundary-mark kernels are obtained by disintegrating
the same physical joint. A raw mark which is not a retained coordinate
is retained through its complete joint conditional kernel, including
its correlations with the other marks. Its value is not replaced by
evaluation at \(P_px\). The physical reverse conditional, the optimizer
reverse conditional, and the actual recognition conditional have their
own roles. Finite Gaussian domination supplies positive smooth versions
of the integrals in this realization. These exact probability
constructions retain all original loop data and all source assignments.
\status{ESTABLISHED}

For an original edge \(e\), its inherited mismatch observable under
either the physical posterior or the restricted optimizer is
\begin{equation}
\mathcal D_e^{p,\mu}(y,S)
=\E_{\mu_\vartheta}
 \left[\tfrac12r_e(X)^2\mid P_pX=y,S\right],
\qquad
\mu_\vartheta\in
\{\boldsymbol\Pi_{\vartheta,\mathsf o},\pi_\vartheta\}.
\label{eq:rg-formation-inherited-edge-statistic}
\end{equation}
This defines two conditional observables, not one common reverse
kernel. Coherence of a retained description refers to their
root-mean-square original residuals
\(\sqrt{2\mathcal D_e^{p,\mu}}\); the local \(d_i\) observables
are defined by the same conditional construction. At \(p=0\) these
reduce to the original configuration residuals. This definition does
not identify a single residual cost with the KL of two full parent
mixtures. \status{DEFINITION}

The root factor of \(\pi_\vartheta^p\) is independent of the assignment
and of \(\vartheta\). On \(M_p\), let \(H_p\) be the Schur complement
of \(H|_{\mathcal K^\perp}\), and let
\(\gamma_{p,\vartheta}(dz\mid y)\) be the cumulative hidden Gaussian
conditional of the quadratic optimizer. Its mean is independent of
\(\vartheta\); its covariance is \(T/\vartheta\) times a fixed
inverse precision. Integrating at fixed \(S\) gives the exact sheet
potential
\begin{equation}
\begin{aligned}
V_{p,S}(y,\vartheta)
 &=\tfrac12\vartheta y^\top H_py+C_p^{\mathrm G}(\vartheta)
       +T\psi_{p,S}(y,\vartheta),\\
\psi_{p,S}(y,\vartheta)
 &=-\log\E_{\gamma_{p,\vartheta}(\cdot\mid y)}
                  e^{-\phi_S},&
0&\le\psi_{p,S}\le B_0,\\
\frac{d}{d\vartheta}C_p^{\mathrm G}(\vartheta)
 &=\frac{Td_p}{2\vartheta}.
\end{aligned}
\label{eq:rg-formation-parent-potential}
\end{equation}
Here \(C_0^{\mathrm G}=0\), \(\psi_{0,S}=\phi_S\); the unchanged
root potential is appended on \(W_p\).
The Gaussian integration scalar \(C_p^{\mathrm G}\) is distinct from
the common objective offset \(C_T\). The potential is an exact
conditional integral, with every generated interaction retained.
\status{ESTABLISHED}

For a density \(q^p\) on the marked parent carrier define
\begin{equation}
F_T^p(q^p,\vartheta)
  =T\KL(q^p\Vert\pi_\vartheta^p)+C_T(\vartheta).
\label{eq:rg-formation-parent-vfe}
\end{equation}
The conditional KL chain rule gives, at one common parameter and one
common time,
\begin{equation}
\begin{aligned}
\Delta_{p:s}(q^p,\vartheta)
 &=\E_{(C_{s\leftarrow p})_\#q^p}
   \KL\left(q^p(dz\mid y,S)
        \middle\Vert\pi_\vartheta^p(dz\mid y,S)\right),\\
F_T^p(q^p,\vartheta)
 &-F_T^s((C_{s\leftarrow p})_\#q^p,\vartheta)
   =T\Delta_{p:s}(q^p,\vartheta).
\end{aligned}
\label{eq:rg-formation-conditional-cost}
\end{equation}
In the displayed conditional, \(z\) consists only of the discarded
continuous normal coordinates; loops and \(S\) remain in the retained
state. Nested conditional losses telescope at that same law and time.
Costs evaluated at different activation times are not summed as one
telescoping objective. \status{ESTABLISHED}

Declare unit continuous mobility on every microscopic coordinate and a
rate-one Poisson proposal clock. At a proposal, redraw the whole
assignment from \(b(x,\vartheta)\) without changing any continuous
coordinate. A proposal may return the current assignment. The microscopic
generator is
\begin{equation}
\mathcal L_\vartheta^0 f(x,S)
=T\Delta_x f-\nabla_x U_{\vartheta,S}\cdot\nabla_xf
 +\sum_{S'}b_{S'}(x,\vartheta)[f(x,S')-f(x,S)].
\label{eq:rg-formation-microscopic-generator}
\end{equation}
The diffusion covariance is \(2TI dt\). This is a continuous-time
inference prescription on global Euclidean space times a finite set,
with integrable decay at infinity and no reflecting or chart boundary.
No discretization is being defined. \status{DEFINITION}

The inherited parent proposal row and generator are
\begin{equation}
\begin{aligned}
R^p_{SS'}(y,\vartheta)
 &=\E_{\pi_\vartheta}
     [b_{S'}(X,\vartheta)\mid P_pX=y,S],\\
\mathcal L_\vartheta^pf(y,S)
 &=T\Delta_yf-\nabla_yV_{p,S}\cdot\nabla_yf
   +\sum_{S'}R^p_{SS'}(y,\vartheta)[f(y,S')-f(y,S)].
\end{aligned}
\label{eq:rg-formation-parent-generator}
\end{equation}
The unchanged root drift is included when \(y\) denotes all of \(W_p\).
Self-proposals are retained, so every row sums to one and every actual
exit rate is at most one. The conservative jump diagonal is
\(R^p_{SS}-1\). These conditional expectations specify the chosen
autonomous parent kinetics; they do not assert exact temporal projection
of the fine process. \status{DEFINITION}

Multiplying the microscopic row by \(\pi_\vartheta(x,S)\) gives a
quantity symmetric in \(S,S'\). Conditional integration therefore
proves detailed balance of \(R^p\) with \(\pi_\vartheta^p\).
Furthermore
\begin{equation}
\begin{aligned}
R^p_{SS'}(y,\vartheta)
 &\ge e^{-B_0}\pi_{S'}(\vartheta)>0,\\
\E_{\pi_\vartheta^p}
 [R^p_{SS'}(Y,\vartheta)\mid P_sY=z,S]
 &=R^s_{SS'}(z,\vartheta).
\end{aligned}
\label{eq:rg-formation-rate-composition}
\end{equation}
The second identity is the conditional-expectation tower. Thus both
normalized laws and optimizer-conditioned rates compose exactly at a
fixed parameter, while comparisons between their independently evolved
histories still require estimates. \status{ESTABLISHED}

\section{Complete model scores and dissipation}
\label{sec:rg-formation-learning}

Differentiate the complete normalized experiment before conditioning on
the observed record. The prior, calibration, and categorical scores
have mean zero and vanishing cross terms: the categorical draws and
calibration are independent of the configuration at fixed
\(\vartheta\), and the remaining conditional kernels have no explicit
\(\vartheta\)-dependence. The complete generative Fisher metric is
therefore
\begin{equation}
M(\vartheta)
=\frac1T+\frac12\Tr[(A_0(\vartheta)^{-1}H)^2]
 +16\pi_{i_0,+}(\vartheta)
       [1-\pi_{i_0,+}(\vartheta)].
\label{eq:rg-formation-model-fisher}
\end{equation}
The quadratic-prior term follows by the variance of its centered
Gaussian quadratic score. This establishes the metric of the original
full experiment; it is not the Fisher metric of a fitted coarse
observation model. \status{ESTABLISHED}

Declare \(m(\vartheta)=TM(\vartheta)\) and the model metric
\(m(\vartheta)/\epsilon\) for the scaled objective, with
\(\epsilon>0\). At every level the actual coupled system is
\begin{equation}
\partial_tq_t^p=(\mathcal L_{\vartheta_p(t)}^p)^*q_t^p,
\qquad
\dot\vartheta_p(t)
  =-\epsilon m(\vartheta_p(t))^{-1}
        g_p(q_t^p,\vartheta_p(t)),
\qquad
g_p=\left.\partial_\vartheta F_T^p\right|_{q^p}.
\label{eq:rg-formation-own-law-dynamics}
\end{equation}
The density is the complete actual law of that level. Its parameter
path is deterministic because this feedback takes an exact-law
expectation, rather than a sampled score. Parents inherit this metric;
neither a new rate multiplier nor a marginal-model Fisher metric is
inserted at activation. \status{DEFINITION}

Expanding \cref{eq:rg-formation-parent-vfe} by
\cref{eq:rg-formation-parent-potential} gives the full covector
\begin{equation}
\begin{aligned}
g_p(q^p,\vartheta)
={}&\vartheta-2+\tfrac12\E_{q^p}y^\top H_py
       +\frac{Td_p}{2\vartheta}
       +T\E_{q^p}\partial_\vartheta\psi_{p,S}(y,\vartheta)\\
 &-\frac T2\Tr(A_0(\vartheta)^{-1}H)
       -T\E_{q^p}\partial_\vartheta\log\pi_S(\vartheta).
\end{aligned}
\label{eq:rg-formation-complete-score}
\end{equation}
The root has no parameter score. The derivative is partial at fixed
actual law, so time derivatives of actual assignment masses are not
extra terms. Prior normalization, hidden determinant, source prior, and
nonlinear conditional normalization are all included. At \(p=0\),
\(d_0=0\) and \(\partial_\vartheta\phi_S=0\); this gives the
microscopic score directly. \status{ESTABLISHED}

For positive smooth densities with finite entropy and moments, write
\(u_S=q^p_S/\pi_\vartheta^p(\cdot,S)\) and
\(c^p_{SS'}(y)=\pi_\vartheta^p(y,S)R^p_{SS'}(y,\vartheta)\).
Integration by parts and the symmetry of this conductance yield
\begin{equation}
\begin{aligned}
\frac{d}{dt}F_T^p
={}&-T^2\sum_S\int q_S^p
       \left|\nabla\log
           \frac{q_S^p}{\pi_{\vartheta_p}^p(\cdot,S)}\right|^2dy\\
 &-\frac T2\int\sum_{S,S'}c^p_{SS'}
       (u_S-u_{S'})(\log u_S-\log u_{S'})dy\\
 &-\epsilon m(\vartheta_p)^{-1}g_p(q^p,\vartheta_p)^2
 \le0.
\end{aligned}
\label{eq:rg-formation-dissipation}
\end{equation}
To obtain the middle line, interchange \(S,S'\) in one copy of the
jump contribution and average. The resulting integrand is nonnegative.
At the microscopic level the same term is
\(-T\int q_X[\KL(\rho_q\Vert b)+\KL(b\Vert\rho_q)]dx\).
The confining moment bounds and cutoff passage in
\Cref{lem:formation-prepared-continuation,lem:formation-conditional-entropy}
justify these identities for the admitted actual solutions. Thus the
declared coupled inference decreases its own complete VFE. This does
not assert that a discarded-conditional KL decreases separately.
\status{ESTABLISHED}

\section{Preparation and successive actual-law activation}
\label{sec:rg-formation-preparation}

The preparation uses a finite nonzero residual margin. Include each
internal edge covector on \(\mathcal N\), and for each bridge \(e\)
include \(r_eP_s|_{\mathcal N}\) for \(s<\ell(e)\), in a list of
\(m_0\) linear functionals. Each is nonzero by residual-row
independence. In a fixed orthonormal basis of \(\mathcal N\), evaluate
these functionals on \((1,t,\ldots,t^{\dim\mathcal N-1})\).
There are at most \(m_0(\dim\mathcal N-1)\) forbidden integers, so
at least one of the integers from zero through that bound is a
simultaneous nonroot. Normalize the resulting vector to obtain
\(v_\star\in\mathcal N\) and a margin
\begin{equation}
m_L=\min_{e,\ s<\ell(e)}
       |r_e(P_sv_\star)|>0.
\label{eq:rg-formation-preparation-margin}
\end{equation}
Internal edges enter this minimum only at \(s=0\). A fixed sufficiently
small neighborhood of \(v_\star\) preserves a margin at least \(m_L/2\).
This is an explicit construction of preparation data, not a time-dependent
merger instruction. \status{ESTABLISHED}

For \(0<\epsilon<1\), fix the parameter family
\begin{equation}
Q_L=0,\qquad Q_s=2Q_{s+1}+20,\qquad
J_s=\epsilon^{-Q_s},\qquad
A^2=\epsilon^{3Q_1+20},\qquad
T=A^2\epsilon^6,\qquad H_*=\frac2\epsilon.
\label{eq:rg-formation-weakest-one-regime}
\end{equation}
The \(Q_s\) are scalar exponents; the orthonormal flag matrices are
\(\mathsf Q_s\). The weakest energy is exactly \(J_L=1\). There is no
common acceleration of every energy, no vanishing source rate, and no
decrease of the standardized nonlinearity \(\delta\).
\status{DEFINITION}

Start \(\vartheta(0)=1\). The root is independent of the non-root
coordinates and assignments, with law
\(\mathcal N_{\mathcal K}(\kappa z_\star/(1+\kappa),
T(1+\kappa)^{-1}I_{\mathcal K})\). On every assignment sheet the
reference non-root law has common mean \(Av_0\), zero loop mean, and
covariance \(T(H|_{\mathcal K^\perp})^{-1}\), where \(v_0\) lies in
the strict-margin neighborhood just constructed. Choose positive
assignment masses near the product prior at \(\vartheta=1\), so
\(|\lambda_{i_0,+}(0)-\operatorname{sigmoid}(-2)|<10^{-2}\).
Allow smooth within-sheet transports of this reference with coupling
\begin{equation}
|\widehat K_1^{1/2}(X_0-X_0^{\mathrm G})|
 <\epsilon\sqrt T
\quad\hbox{almost surely},
\label{eq:rg-formation-preparation-transport}
\end{equation}
and bounded log-density ratio, retaining the independent root factor.
Small smooth compactly supported diffeomorphisms in standardized
non-root coordinates supply such transports. The resulting neighborhood
of positive smooth correlated initial laws is understood in this weighted
transport and bounded-log-density topology. Its size may depend on
\(\epsilon\) and finite \(L\). These are preparation hypotheses;
the later actual law is not constrained to remain Gaussian or to
factor across assignment sheets. \status{HYPOTHESIS}

Choose one finite entrance constant \(k_L\) as in the theorem below,
and set
\begin{equation}
t_s=\frac{k_L\log(1/\epsilon)}{J_s},
\qquad I_s=[t_s,2t_s],\qquad
\tau_0=0,\quad\tau_s=2t_s.
\label{eq:rg-formation-activation-times}
\end{equation}
At \(\tau_s\), push the complete actual level-\((s-1)\) law through
\(C_{s\leftarrow s-1}\), copy its actual parameter, and thereafter
evolve the new law by
\cref{eq:rg-formation-own-law-dynamics}. An independently continued
source process retains its own law and parameter after that time.
These deterministic times choose when to use a reduced description;
they do not change the microscopic graph, energy hierarchy, or
inference rule. \status{DEFINITION}

\section{The finite-depth formation theorem}
\label{sec:rg-formation-theorem}

\theoremheading{Adaptive formation with weakest energy one}
{thm:finite-depth-adaptive-formation}
For every prescribed finite \(L\ge2\), fix the graph, evaluator,
source support, priors, observations, and microscopic kinetics above.
There exist positive finite constants, depending on these fixed data
and \(L\), and \(\epsilon_L>0\), such that for every
\(0<\epsilon<\epsilon_L\) the regime
\cref{eq:rg-formation-weakest-one-regime} and the declared initial-law
neighborhood admit all \(L\) successive actual-law activations.
Choose
\begin{equation}
\begin{aligned}
k_L&>k_{\mathrm{geom},L},&
\tfrac12c_{\mathrm{early},L}k_L&>Q_1+20,\\
c_{\mathrm{joint},L}k_L&>Q_1+8,&
c_{\mathrm{fast},L}k_L&>Q_1/2+12,
\end{aligned}
\label{eq:rg-formation-common-entrance}
\end{equation}
where the positive restoring, joint-entropy, and fast-mixing constants
are those constructed in
\Cref{lem:formation-coherence,lem:formation-conditional-entropy,lem:formation-history-coupling}.
The finite geometric entrance constant is constructed from
\cref{eq:rg-formation-weighted-flag,eq:rg-formation-preparation-margin}
in the coherence lemma. The following conclusions are established
for the resulting actual processes. \status{ESTABLISHED}

The coupled diffusion, finite-state proposals, and exact-law parameter
equation have a unique nonexplosive solution through \(H_*\).
Every activated parent retains positive smooth densities and satisfies
\begin{equation}
\begin{aligned}
1&\le\vartheta_p(t)\le21/10,&
|\dot\vartheta_p(t)|&\le C_L\epsilon,\\
|\vartheta_p(t)-(2-e^{-\epsilon t})|
 &\le C_Lr_\epsilon,&
r_\epsilon&=A^2J_1+
       T[1+\log(1/\epsilon)]^2.
\end{aligned}
\label{eq:rg-formation-parameter-bound}
\end{equation}
Here the time in the comparison is absolute time, including after
activation. The reference \(2-e^{-\epsilon t}\) is an estimate
deduced from the complete own-law equation, not a prescribed update.
\status{ESTABLISHED}

Choose a fixed residual tolerance \(\eta_LA\) smaller than the
strict prepared separation margin. Except on an event of probability
at most
\begin{equation}
C_L(1+H_*)\exp(-c_L/\epsilon^6),
\label{eq:rg-formation-coherence-probability}
\end{equation}
the microscopic process and its sequential parent descriptions have
the successive coherent components
\begin{equation}
3\cdot2^{L-1}\longrightarrow2^{L-1}
\longrightarrow2^{L-2}\longrightarrow\cdots\longrightarrow1.
\label{eq:rg-formation-agent-counts}
\end{equation}
On \(I_s\), edges already assigned to levels at most \(s\), and the
within-node model discrepancies, satisfy the coherence tolerance;
future bridges remain separated by a fixed positive multiple of \(A\).
The retained description uses
\cref{eq:rg-formation-inherited-edge-statistic} and its local
model-discrepancy counterpart. The final coherent state persists through the
finite horizon on this event. All loop coordinates remain random and
retained; no identity-holonomy constraint is imposed.
\status{ESTABLISHED}

On each active interval \(I_s\), let
\(\Delta_{s-1:s}\) be the actual discarded conditional loss of
\cref{eq:rg-formation-conditional-cost}. At the same current law
and parameter,
\begin{equation}
\begin{aligned}
\Delta_{s-1:s}
 &\le
 \begin{cases}
C_L\epsilon^{14},&s<L,\\
C_L\epsilon^2,&s=L,
\end{cases}\\
|g_{s-1}(q,\vartheta)
 -g_s((C_{s\leftarrow s-1})_\#q,\vartheta)|
 &\le
 \begin{cases}
C_LT\epsilon^7,&s<L,\\
C_LT\epsilon,&s=L.
\end{cases}
\end{aligned}
\label{eq:rg-formation-conditional-score-bounds}
\end{equation}
Multiplication of the first bound by \(T\) gives the discarded
restricted VFE. These are conditional-entropy and complete-score
estimates, independently derived from the actual evolving laws.
\status{ESTABLISHED}

Let \(\lambda^p_{i_0,+}(t)\) be the actual distinguished receiver
source mass in a fine, parent, or sequential process. For each actual
activation time from which that process is continued,
\begin{equation}
\lambda^p_{i_0,+}(H_*)-
  \lambda^p_{i_0,+}(\tau_s)>\tfrac12.
\label{eq:rg-formation-source-movement}
\end{equation}
Thus the actual binary recognition-source row changes by total
variation greater than one half. Moreover, between
\(\vartheta=1\) and the final learned parameter, both the first hidden
physical conditional and the first hidden optimizing conditional change
by KL greater than \(1/200\), at every fixed retained input and
assignment. These compare their own normalized kernels, not one kernel
substituted for another. \status{ESTABLISHED}

For any source process and its inherited parent begun from a common
actual marked activation sample at time \(\tau\), there exists a
coupling preserving that sample for which
\begin{equation}
\begin{aligned}
\Pr\{S_t\ne\overline S_t
       \text{ for some }t\in[\tau,H_*]\}
 &\le C_L\epsilon^6,\\
\Pr\left\{\sup_{\tau\le t\le H_*}
       |P_sX_t-\overline X_t|>C_LA\epsilon^2\right\}
 &\le C_L(1+H_*)e^{-c_L/\epsilon^2}.
\end{aligned}
\label{eq:rg-formation-two-path-bounds}
\end{equation}
The first comparison is of the complete assignment path, with every
receiver retained. The second is a uniform continuous-coordinate
displacement bound under a compatible coupling. A finite number of
successive replacements changes only the constants. No total variation
convergence of continuous path laws or exact Markov lumpability is
asserted. \status{ESTABLISHED}

\paragraph{Proof.}
The exact physical normalization, restricted scalar, and marked
composition have been calculated in
\cref{eq:rg-formation-full-joint,eq:rg-formation-evidence-scalar,eq:rg-formation-conditional-cost,eq:rg-formation-rate-composition}.
\Cref{lem:formation-parent-score} proves the uniform differentiated
parent bounds, including their hidden determinant and source terms.
\Cref{lem:formation-prepared-continuation} constructs the coupled
solution, uses stopped moments to keep the parameter in its positive
interval, and propagates the admitted moments through actual
activations. \Cref{lem:formation-coherence} proves the strict
formation margins from the weighted flag.
\Cref{lem:formation-parameter-comparison} gives
\cref{eq:rg-formation-parameter-bound} for separately learning
parents. These arguments use neither conditional equilibration nor
history averaging.

For early interfaces, \Cref{lem:formation-conditional-entropy}
keeps the finite positive jump contribution in the conditional balance.
The contrast identity \(J_{s+1}^2/J_s=\epsilon^{20}\) and the
entrance condition in \cref{eq:rg-formation-common-entrance} give
the first early bound in
\cref{eq:rg-formation-conditional-score-bounds}. At the terminal
interface the same lemma instead uses joint diffusion-plus-jump
entropy contraction after logarithmic time. Conditional KL is bounded
by that full entropy through the chain rule. The complete normalized
hidden score then bounds the score difference by
\(C_LT(\sqrt\Delta+\Delta)\), giving the second line of
\cref{eq:rg-formation-conditional-score-bounds}.

\Cref{lem:formation-source-envelope} bounds the actual source mass
at every activation, including the final logarithmic activation, and
integrates its positive-rate equation on the last ten units of the
horizon to prove \cref{eq:rg-formation-source-movement}. The same
parameter continuation supplies the hidden-kernel comparison proved
there. \Cref{lem:formation-fast-marginal} integrates later triangular
residuals exactly, leaving a common source-relevant fast optimizer
marginal at every level. \Cref{lem:formation-history-coupling}
then compares conditional proposal-prefix kernels under the actual
evolving environment. Its error is \(C_L\epsilon^6\), with the
common actual activation information retained. Finally
\Cref{lem:formation-continuous-path} supplies the separate
continuous bound on a coupling compatible with those histories.
Each lemma uses the same preparation, parameter regime, and finite
choice \cref{eq:rg-formation-common-entrance}; the final window is
of order \(\log(1/\epsilon)\) and lies before \(H_*=2/\epsilon\)
for small \(\epsilon\). This proves the conjunction. \(\square\)

\section{A two-triangle instance and the remaining scope}
\label{sec:rg-formation-example-scope}

For \(L=2\), take triangles \((1,2,3)\) and \((4,5,6)\), with a
single bridge \((3,4)\). The retained carriers are
\begin{equation}
\R^{19}\times\mathcal J
\longrightarrow\R^9\times\mathcal J
\longrightarrow\R^8\times\mathcal J,
\qquad |\mathcal J|=64,
\label{eq:rg-formation-two-triangle-carrier}
\end{equation}
and the coherent agent counts are \(6\to2\to1\).
Both loops survive both maps. The exact regime becomes
\begin{equation}
J_1=\epsilon^{-20},\qquad J_2=1,\qquad
A=\epsilon^{40},\qquad T=\epsilon^{86}.
\label{eq:rg-formation-two-triangle-regime}
\end{equation}
The first window has order
\(\epsilon^{20}\log(1/\epsilon)\), and the second has order
\(\log(1/\epsilon)\). The source and model changes occur over the
longer horizon \(2/\epsilon\). These are exact substitutions into
the finite-\(L\) theorem, not numerical evidence for it.
\status{ESTABLISHED}

The source-dependent term is nonquadratic in the declared
configuration coordinates. For a direction with selected internal
residual slope \(\lambda\ne0\),
\(f(\lambda t)=\lambda^2t^2/2-\lambda^4t^4/8+O(t^6)\).
Every assignment selects an internal edge with nonzero first-normal
projection, so its hidden log density has a nonzero fourth derivative
in some first-normal direction. Exact marginalization retains its full
conditional integral and the negative covariance term in its Hessian.
Thus neither the physical nor optimizing first hidden conditional is
silently replaced by a Gaussian conditional. This is an established
coordinate-specific nonlinearity statement. \status{ESTABLISHED}

The theorem keeps the supplied energy hierarchy and its internal source
support. It does not prove that source learning creates those contrasts.
Allowing sources on different bridge levels would require a new
fast-marginal and evolving-environment estimate. Uniform control as
\(L\to\infty\), a nondegenerate intrinsic partition selector,
permanent membership under arbitrary later forcing, and autonomous
agency remain separate open construction problems. The finite
constants and admissible preparation neighborhoods need not be uniform
in depth. \status{OPEN}

The parameter \(t\) in the generators is the declared inference time.
This chapter declines to identify it with physical time or with Fisher
length, and it does not claim computational savings or a
depth-independent resource bound. Exact common-channel probability
composition, approximate continuous trajectories, and complete
assignment-history comparison retain the separate meanings used in
\Cref{ch:agent-network-rg,ch:relational-inference}.
\status{NOT-CLAIMED}
